% =====================================================================
% sections/discussion.tex  --  input by report.tex.
%
% Word budget: ~800 words (counted toward the 3000-word limit).
%
% Spec: do NOT list the discussion questions and answer them in turn -- weave.
% Required coverage, woven: (a) further work; (b) how adaptive the system is,
% per the intro's definition; (c) connection to the wider adaptive-systems
% area; plus at least one of hypothesis-support / new-hypothesis.
%
% Builds to the two main points:
%   MP1: genetic assimilation has NOT occurred (frozen-HP test).
%   MP2: the Dev only paradox -- best-performing yet most HP-dependent.
%
% "Consistent with" discipline preserved from methods_log 11.5: the 8b/8c
% conjunction is consistent with Williams's substrate claim, not a proof; the
% Both mechanism is the most parsimonious reading, not a demonstrated cause.
% =====================================================================

\section{Discussion}
\label{sec:discussion}

The central question of this project, framed through the Baldwin effect, was
whether behaviour evolved with homeostatic plasticity becomes genetically
assimilated --- encoded innately in the genotype --- or remains dependent on the
plasticity that produced it \parencite{baldwin1896, hintonnowlan1987}. The
frozen-HP test answers this directly, and the answer is that assimilation has
not occurred in any HP condition. Across all three, removing HP by adiabatic
elimination drops fitness far beyond the measurement-noise baseline, and in a
substantial fraction of individuals --- between three and five in ten depending
on condition --- it collapses the behaviour entirely. The within-lifetime
adaptation is not a refinement layered on top of a competent genotype; in these
individuals it is constitutive of the behaviour. This is the clearest result of
the study, and it is a negative result with respect to the Baldwin prediction:
selection under HP did not, over $200$ generations, produce genotypes that no
longer need it.

The viable-range diagnostics make the same point at the level of neural
dynamics. Every HP condition holds neurons in the viable range while HP is
active and loses that property the moment HP is frozen, falling to or below the
non-plastic baseline. The same intervention collapses both viable-range
operation and fitness, which is \emph{consistent with} Williams's claim that
viable-range operation is functionally important
\parencite{williams2006} without proving it, since fitness could be sensitive
to the frozen parameters for reasons the viable-range fraction does not capture.
Stolting et al.\ offer a candidate mechanism \parencite{stolting2023}:
HP-during-behaviour may evolve dynamics in the joint network-and-parameter state
space that cannot be reproduced once the parameters are held fixed. The large
drops are consistent with this but do not isolate it from simpler explanations.

The second main finding is an apparent paradox. \emph{Dev only} is both the
best-performing condition and the most HP-dependent, with the highest rate of
complete collapse on freezing. These are not in tension once the selection
regime is considered. Because HP always ran before evaluation in \emph{Dev
only}, every genotype the GA ever scored had already been transformed by
development; the GA could not distinguish a genotype that works because its raw
parameters are good from one that works only after HP has reshaped it. Both look
identical to the fitness function, and selection therefore had no purchase on
innate competence at all. The high fitness and the high HP-dependence are two
faces of the same fact. This is the Baldwin setup operating without its second
phase: lifetime adaptation guided the search, but because acquired changes are
not inherited and there was never selective pressure to internalise them, no
assimilation followed. It also resolves the puzzle of why \emph{Both}, which
runs the identical developmental phase, fails to inherit \emph{Dev only}'s
advantage: leaving HP active during behaviour appears to let ongoing plasticity
work against the developmentally established configuration rather than build on
it. That reading is the most parsimonious account of the data rather than a
directly demonstrated cause.

One confound should be stated plainly. The conditions vary both the presence of
a developmental phase and the presence of behaviour-phase HP, but do not include
a developmental phase with HP disabled, so the developmental benefit cannot be
fully separated from the effect of simply letting the network settle for $6000$
unloaded timesteps before evaluation. The evidence points against settling as
the primary driver --- the substrate check shows a genuine parameter shift, and
the collapsed individuals have raw genotypes whose dynamics no amount of
settling would rescue --- but the question is not closed for the moderate
effects.

The search-dynamics analysis reframes all of this at the level of the larger
adaptive system. The four conditions converge on the same early timescale and
differ only in where they settle, which locates HP's effect precisely: it does
not change the shape of the fitness landscape, which is fixed by the evaluation
procedure, but changes the effective genotype-to-phenotype mapping the GA
selects across. \emph{Dev only} reaches a higher basin not by searching longer
but by handing the GA a better-developed phenotype to score. Read this way, the
one respect in which the replication departs from Williams --- that online HP
was not significantly worse than no HP, where he reports it as worse --- is
plausibly an effect of the better-conditioned tournament selection used here,
which is less prone than his roulette-plus-elitism scheme to amplify small
fitness differences into convergence on a poor basin. The point is offered as
the most likely reading rather than a tested claim, since the two GAs were not
run side by side.

Returning to the working definition of adaptivity from the introduction --- a
system that maintains viability under perturbation by acting across multiple
timescales \parencite{ashby1960, dipaolo2000} --- the HP conditions are more
adaptive than the non-plastic one in a precise and somewhat uncomfortable sense.
HP is genuine cybernetic negative feedback, regulating firing rates toward a
viable set-point much as Ashby's ultrastable system reorganises to stay within
bounds. But the frozen-HP results show that for many evolved individuals this
regulation is load-bearing rather than corrective: the agent does not merely
perform better with HP, it requires HP to function. The adaptivity has not been
absorbed into the structure; it remains an active process the behaviour leans
on.

Several directions follow. The cleanest is the missing fifth condition --- a
developmental phase with HP disabled --- which would resolve the settling
confound. The discrimination task Williams found hardest would test a clear
prediction: if online HP is constitutive on the easy task, the dependence should
deepen on the harder one. Running Williams's original roulette-plus-elitism GA
alongside the present tournament GA would separate selection effects from HP
effects, the one methodological difference this replication retains. Finally,
the assimilation question itself could be pursued by extending evolution well
beyond $200$ generations, or by adding an explicit cost to plasticity, to ask
whether assimilation occurs at all under stronger or more prolonged selective
pressure \parencite{mayley1996}.
